% \section{Các công trình nghiên cứu liên quan}

\begin{frame}{Những hướng tiếp cận chính}
    % Có 3 hướng tiếp cận chính:
    \begin{itemize}
        \item Học trên cấu trúc cơ thể.
        \item Học trên ngoại hình/hình dáng.
        \item Học trên đa phương thức.
    \end{itemize}

\end{frame}

% \begin{frame}{Học trên cấu trúc cơ thể}
%     chèn ảnh minh họa kiểu dữ liệu, các tập thường dùng, các mô hình liên quan.
%     \begin{itemize}
%         \item Tập dữ liệu: OU-MVLP(Pose/Skeleton), Gait3D, GREW.
%         \item Các phương pháp: GaitGraph, PoseGait, Gait-TR, GPGait, SkeletonGait.
%     \end{itemize}
% \end{frame}

% TODO: hãy bổ sung tham chiếu/tham khảo cho từng tập dữ liệu, phương pháp được liệt kê. Phải tìm kiếm từ tác cả nguồn trong thư mục utils
\begin{frame}{Học trên cấu trúc cơ thể}
    \begin{itemize}
        \item Đầu vào: Cấu trúc cơ thể được ước lượng (khung xương, mô hình cơ thể 3D.% (SMPL \cite{loper2015smpl}))
        \item Nguyên lý: Tham số hóa dáng đi dưới dạng các vec-tơ tọa độ/đồ thị xương, loại bỏ hầu hết các thông tin về màu sắc hay kết cấu bề mặt.
        \item Ưu điểm:
              \begin{itemize}
                  \item Đặc trưng \textit{"sạch"}, tức là không chứa thông tin gây nhiễu như trang phục hoặc phụ kiện.
                  \item Độ hiệu quả bất biến với sự thay đổi về trang phục hay vật mang theo, do được huấn luyện dựa trên chuyển động của khớp xương.
              \end{itemize}
        \item Nhược điểm
              \begin{itemize}
                  \item Phụ thuộc nhiều và độ chính xác của thuật toán ước lượng tư thế.
                  \item Độ hiệu quả kém với các tập dữ liệu thế giới thực, vì bỏ qua thông tin về dáng người thô.
              \end{itemize}
        % \item Tập dữ liệu: OU-MVLP (Pose) \cite{takemura2018multi}, Gait3D \cite{zheng2022gait}, GREW \cite{zhu2021gait}.
        % \item Các phương pháp: GaitGraph \cite{teepe2021gaitgraph}, PoseGait \cite{liao2020model}, Gait-TR \cite{zhang2021gait}, GPGait \cite{fu2023gpgait}, SkeletonGait \cite{fan2024skeletongait}.
    \end{itemize}
\end{frame}

\begin{frame}{Học trên ngoại hình/hình dáng}
    \begin{itemize}
        \item Đầu vào: chuỗi hình nhị phân của bóng/dáng người, hoặc ảnh RGB. Hình dáng được xem xét trong nhiều tình huống (ví dụ: trang phục bình thường, nhiều trang phục, mang túi xách...).
        \item Nguyên lý: Học trực tiếp từ chuỗi dáng đi, bằng mạng học sâu, để trích xuất đặc trưng hình dáng (không gian) và chuyển động (thời gian).
        \item Ưu điểm:
              \begin{itemize}
                  \item Tận dụng thông tin hình dáng cơ thể đa dạng, tăng khả năng phân biệt danh tính.
                  \item Nhận được nhiều sự quan tâm, và đạt hiệu suất cao trên các tập dữ liệu chuẩn trong phòng thí nghiệm.
              \end{itemize}
        \item Nhược điểm:
              \begin{itemize}
                  \item Dễ bị ảnh hưởng bởi các yếu tố gây nhiễu trong thực tế (vật che khuất phức tạp, hậu cảnh, sự thay đổi ánh sáng, góc nhìn.)
                  \item Hiệu suất giảm khoảng hơn 40\% trong môi trường thực tế, so với phòng thí nghiệm.
              \end{itemize}
        % \item Tập dữ liệu: CASIA-B \cite{yu2006framework}, OU-MVLP \cite{takemura2018multi}, Gait3D \cite{zheng2022gait}, GREW \cite{zhu2021gait}, CCPG.
        % \item Phương pháp tiêu biểu: GaitSet \cite{chao2019gaitset}, GaitGL \cite{lin2021gait}, GaitBase, GaitPart \cite{fan2020gaitpart}, 3D Local \cite{huang20213d}.
    \end{itemize}
\end{frame}

% \begin{frame}{Học trên ngoại hình/hình dáng}
%     \begin{itemize}
%         \item Tập dữ liệu: CASIA-B, OU-MVLP, Gait3D, GREW, CCPG
%         \item Phương pháp tiêu biểu: GaitSet, GaitGL, DeepGaitV2, GaitBase, DANet, GaitSSB, QAGait.
%     \end{itemize}
% \end{frame}

\begin{frame}{Học trên đa phương thức}
    % Có 3 hướng tiếp cận chính:
    \begin{itemize}
        \item Nguyên lý: kết hợp cùng lúc nhiều đầu vào (hình bóng, khung xương,...) nhằm tạo ra đặc trưng dáng đi toàn vẹn hơn.
        \item Ưu điểm
              \begin{itemize}
                  \item Tận dụng sự bổ trợ lẫn nhau giữa các loại dữ liệu (dữ liệu hình bóng cung cấp thông tin hình dáng và chuyển động; dữ liệu khung xương cung cấp cấu trúc chuyển động bất biến với nhiễu).
                  \item Thường đạt độ chính xác cao hơn và ổn định hơn trong các tình huống khó khăn.
              \end{itemize}
        \item Nhược điểm
              \begin{itemize}
                  \item Độ phức tạp tính toán cao do xử lý nhiều luồng dữ liệu.
                  \item Yêu cầu các luồng dữ liệu phải đồng bộ và có chất lượng cao từ nhiều nguồn cảm biến hoặc thuật toán tiền xử lý khác nhau.
              \end{itemize}
        % \item Tập dữ liệu: Gait3D \cite{zheng2022gait}, SUSTech1K, CCPG.
        % \item Phương pháp tiêu biểu: SMPLGait \cite{zheng2022gait}, BiFusion \cite{hou2020bifusion}, ParsingGait \cite{zheng2023parsing}, SkeletonGait++ \cite{fan2022skeletongait}, GaitRef \cite{zhu2023gaitref}, HybridGait \cite{dong2024hybridgait}.
    \end{itemize}
\end{frame}

% \begin{frame}{Học trên ngoại hình/hình dáng}
%     \begin{itemize}
%         \item Tập dữ liệu: CASIA-B, OU-MVLP, Gait3D, GREW, CCPG
%         \item Phương pháp tiêu biểu: GaitSet, GaitGL, DeepGaitV2, GaitBase, DANet, GaitSSB, QAGait.
%     \end{itemize}
% \end{frame}